\documentclass[11pt]{article}
\usepackage{graphicx} % Required for inserting images
\usepackage{hyperref}
\title{ECE567 Phase 1 Report}
\author{
Group 15 \\
Yunqi Zhao, Yiming Feng, Zhongqi Huang \\
Hyunjae Kim, Vansh Thakur, Daniel Dosti
}
\date{March 2026}

\begin{document}

\maketitle


\section*{Project 2: Online Goal-Conditioned Reinforcement Learning}
\subsection*{Introduction}

We explore online goal-conditioned reinforcement learning (GCRL) using the JaxGCRL benchmark. In Phase 1, our objective is to reproduce several baseline algorithms, including Contrastive Reinforcement Learning (CRL), goal-conditioned SAC, PPO, and TD3.
\newline
\noindent
We evaluate these methods across a variety of tasks provided within JaxGCRL, spanning simpler environments (e.g., reacher, cheetah), locomotion tasks (e.g., ant, humanoid), navigation-based maze environments, and manipulation tasks (e.g., arm based control). This diversity allows us to assess how different algorithms perform across varying levels of difficulty and task structure.
\newline
\noindent
We selected following environments for comparison.
\newline
\noindent
\texttt{ant}, \texttt{ant\_push},\texttt{ant\_random\_start}, \texttt{humanoid},\texttt{cheetah}
\subsection*{Reproducibility}
Reproducibility of JaxGCRL is sensitive to computing environment based on Phase 1 experience. For instance, the basic running script made development environment errors without a sophisticated CUDA setup. Furthermore, JaxGCRL requires significantly high sample complexity. In order to reproduce the exact same result in \cite{bortkiewicz2025acceleratinggoalconditionedrlalgorithms}, it requires more than 100 million of environment steps. Although we ran the JAXGCRL without modification, due to lack of computation resource, we were not able to achieve the exactly same result shown in \cite{bortkiewicz2025acceleratinggoalconditionedrlalgorithms}. For instance, for \texttt{ant\_push}, it recorded \verb|assert total_steps >= config.total_env_steps|\verb|AssertionError| with a default running script. Furthermore, cheetah environment had crash with \verb|KeyError: 'distance_from_origin'|
% \begin{table}[h]
% \centering
% \small
% \begin{tabular}{l|ccccc}
% \hline
% Env & TD3 & PPO & SAC & CRL \\
% \hline
% ant & 1 & 1 & 1 &421.64453\\
% ant push & 1 & 1 & 1 &263.96875
% \\
% ant random start & 1 & 1 & 1 &1
% \\

% cheetah & 1 & 0 & 1 & N/A
% \\

% humanoid & 1 & 0 & 17.96094 &28.55078
% \\

% \hline
% \end{tabular}
% \caption{Comparison between TD3, PPO, SAC, and CRL with final episode reward}
% \label{tab:reproduce}
% \end{table}


\noindent{TD3 was evaluated on several locomotion environments within the JaxGCRL benchmark, including \texttt{ant}, \texttt{ant\_random\_start}, and \texttt{humanoid}, with all experiments run for approximately 50 million environment steps. Across all environments, TD3 failed to learn a meaningful goal-conditioned policy, achieving zero episode reward and zero success rate in both \texttt{ant} and \texttt{ant\_random\_start}, while primarily accumulating survival reward without making progress toward the goal. Similarly, in the \texttt{humanoid} environment, TD3 did not demonstrate consistent goal-reaching behavior, with only negligible success in easier cases. These results are consistent with prior findings from the JaxGCRL benchmark, where TD3 performs poorly in sparse-reward, goal-conditioned settings due to limited exploration and the absence of mechanisms such as hindsight relabeling or contrastive learning. As rewards are only provided upon reaching the goal, TD3 struggles to discover successful trajectories and instead converges to policies that maximize survival rather than task completion, highlighting its limitations in online goal-conditioned reinforcement learning.}
\noindent{The issue continues with PPO, where sparse rewards continue to negatively affect performance. The poor performance across a host of environemnts confirms the results of the paper, which reports poor PPO performance as well}
\subsection*{GitHub Repo}
We included plot results in the following GitHub repository link. The folder name is "ECE567Reproducibility".
\href{https://github.com/danpro1011/ECE-567-P2-Replication}{Team 15 Project2 GitHub Repository}

\newpage
\section*{Project 3: Offline Goal-Conditioned Reinforcement Learning}
\subsection*{Introduction}

This report summarizes our reproduction of offline goal-conditioned reinforcement
learning (GCRL) baselines from OGBench (Park et al., ICLR 2025). OGBench is a
benchmark spanning 85 dataset--task combinations across locomotion mazes, robotic
manipulation, and pixel-based drawing environments. We reproduced 5 algorithms
across 9 selected datasets (3 seeds each), for a total of 135 training runs.

We report overall\_success (the primary OGBench metric) for each algorithm and
dataset, and discuss reproducibility observations.

\paragraph{Environments}

We selected 9 datasets covering three categories:

\textbf{Locomotion:}
antmaze-large-navigate-v0,
antmaze-large-stitch-v0,
antmaze-teleport-navigate-v0,
humanoidmaze-medium-navigate-v0.

\textbf{Manipulation:}
cube-double-play-v0,
scene-play-v0,
puzzle-3x3-play-v0.

\textbf{Powderworld:}
powderworld-medium-play-v0,
powderworld-hard-play-v0.

Locomotion and manipulation environments use MuJoCo/dm\_control. Powderworld is a
pixel-based environment requiring a convolutional encoder (impala\_small) and
discrete actions.

\paragraph{Baselines}

We reproduced five algorithms:

CRL, HIQL, QRL, GCIQL, GCIVL.

All experiments used official hyperparameters without modification.

\paragraph{Results Summary}

Results follow consistent patterns:

\begin{itemize}
\item HIQL dominates locomotion tasks (e.g., antmaze and humanoidmaze)
\item GCIQL performs best in manipulation tasks
\item GCIVL is the only effective method in Powderworld
\item QRL performs well in locomotion but fails in manipulation
\item CRL is competitive in locomotion but weak elsewhere
\end{itemize}


\subsection*{Reproducibility}

Experiments were conducted on the University of Michigan Great Lakes HPC cluster
using SLURM and A100 GPUs.

We used the official OGBench codebase and hyperparameters without modification.
Minor adjustments were made for cluster compatibility (e.g., dataset paths and
logging flags).

Our reproduced results closely match those reported in the original paper, with
differences generally within one standard deviation. These discrepancies are
likely due to differences in random seeds and system environments (e.g., CUDA
versions). The results for locomotion tasks are shown below and the detailed comparison can be seen in our Github repository.

\begin{table}[h]
\centering
\small
\begin{tabular}{l|l|ccccc}
\hline
Task & Source & GCIVL & GCIQL & QRL & CRL & HIQL \\
\hline

\multicolumn{7}{c}{\textbf{Locomotion}} \\
\hline
antmaze-large-navigate 
& Paper & 0.16 & 0.34 & 0.75 & 0.83 & 0.91 \\
& Ours  & 0.149 & 0.327 & 0.807 & 0.865 & 0.889 \\

antmaze-large-stitch 
& Paper & 0.18 & 0.07 & 0.18 & 0.11 & 0.67 \\
& Ours  & 0.188 & 0.111 & 0.212 & 0.187 & 0.677 \\

antmaze-teleport-navigate 
& Paper & 0.39 & 0.35 & 0.35 & 0.53 & 0.42 \\
& Ours  & 0.391 & 0.303 & 0.327 & 0.539 & 0.404 \\

humanoidmaze-medium-navigate 
& Paper & 0.24 & 0.27 & 0.21 & 0.60 & 0.89 \\
& Ours  & 0.251 & 0.309 & 0.193 & 0.576 & 0.907 \\

\hline
\end{tabular}
\caption{Comparison between reported OGBench results and our reproduced results.}
\label{tab:reproduce}
\end{table}

\paragraph{Observations}

\begin{itemize}
\item HIQL shows higher variance in harder tasks (e.g., antmaze-stitch)
\item GCIQL exhibits instability in pixel-based environments
\item Powderworld tasks remain challenging for all methods
\end{itemize}

\paragraph{Limitations}

\begin{itemize}
\item Only 3 seeds were used (paper uses 5)
\item Not all datasets were reproduced
\item Some variance remains unexplained due to hardware differences
\end{itemize}

\paragraph{Conclusion}

Our reproduction confirms that no single method dominates across all environments.
HIQL is best for locomotion, while IQL-based methods (GCIQL/GCIVL) are essential
for manipulation and pixel-based tasks.





\bibliographystyle{ieeetr}   
\bibliography{ref}
\end{document}
